\chapter{Related Works}


This chapter reviews existing solutions for ML experiment caching and identifies 
the specific gap this thesis addresses: the lack of validated best practices 
for caching in Leave-One-Subject-Out (LOSO) cross-validation scenarios.

\section{Academic Research on ML Caching}

Recent research explores adaptive caching strategies for ML workloads. Krishna 
(2025) surveys ML-aided caching techniques including LSTM-based cache usage 
prediction and reinforcement learning for adaptive cache policies~\cite{krishna2025cache}. 
These approaches focus on \textit{learned} caching decisions using predictive 
models.

In the context of reproducibility, Schubotz et al. (2022) demonstrate how 
caching can accelerate data science experiments while maintaining FAIR 
principles~\cite{schubotz2022caching}. Their work emphasizes content-addressable 
storage for automatic result reuse, similar to the fingerprinting approach 
used in this thesis.

However, existing academic work on ML caching does not specifically address 
the requirements of cross-validation experiments, where cache validity depends 
on explicit encoding of data splits. The challenge of fold-aware cache key 
generation remains unaddressed in the literature.

\section{Commercial MLOps Tools}

Current MLOps platforms provide general experiment tracking and pipeline 
caching but lack native support for cross-validation-specific scenarios:

\begin{itemize}
\item \textbf{MLflow:} Provides experiment tracking, model registry, and 
      basic result caching~\cite{zaharia2018mlflow}. Pipeline caching requires 
      the MLflow Recipes module, which caches individual steps but does not 
      natively encode fold assignments in cache keys.

\item \textbf{DVC:} Uses MD5 hashes of pipeline inputs and code to enable 
      pipeline stage reuse~\cite{dvc2022inproceedings}. While DVC can cache 
      cross-validation results, it requires manual configuration of the pipeline 
      definition to parameterize fold indices—a process not documented in 
      standard DVC workflows.

\item \textbf{ZenML:} Implements hash-based step caching with automatic 
      invalidation when inputs or code change. However, ZenML's caching 
      abstractions do not provide explicit support for encoding held-out 
      subject identifiers, requiring users to implement custom materializers.

\item \textbf{Neptune AI and Weights \& Biases:} Excel at experiment logging, 
      metric visualization, and collaborative features~\cite{neptune2024docs,wandb2024docs}, 
      but do not provide computational caching mechanisms—only metadata tracking.
\end{itemize}

\section{Research Gap Analysis}

Table~\ref{tab:related_works_comparison} compares caching capabilities of 
existing tools relevant to this thesis.


\begin{center}
\colorbox{red!20}{%
\begin{minipage}{0.9\textwidth}
\textcolor{red}{\textbf{TODO:}} Confirm on latest tool versions + format table new
\end{minipage}
}
\end{center}


\begin{table}[ht]
\centering
\begin{tabular}{|p{2.5cm}|p{2.5cm}|p{2.5cm}|p{5cm}|}
\hline
\textbf{Tool/Framework} & \textbf{Caching Strategy} & \textbf{Fingerprinting Method} & \textbf{LOSO Support} \\
\hline
MLflow & Step-level (optional) & Parameter hashing & Manual fold parameterization \\
\hline
DVC & Pipeline stage & MD5 (code + inputs) & Manual pipeline params \\
\hline
ZenML & Step-level & SHA-256 (inputs + code + env) & Custom materializers \\
\hline
Neptune AI / W\&B & None (tracking only) & N/A & N/A \\
\hline
\textbf{This Thesis} & \textbf{Hierarchical (features + models)} & \textbf{SHA-256 with subject ID} & \textbf{Native fold-aware keys} \\
\hline
\end{tabular}
\caption{Comparison of caching capabilities in MLOps tools. ``LOSO Support'' 
indicates whether the tool natively encodes fold assignments in cache keys 
or requires manual configuration.}
\label{tab:related_works_comparison}
\end{table}

The core gap is the lack of validated, documented patterns for 
implementing fold-aware caching in cross-validation workflows. While tools 
like DVC and ZenML provide general pipeline caching infrastructure, adapting 
them for LOSO requires custom logic that:
\begin{enumerate}
\item Encodes the held-out subject identifier in cache keys to prevent 
      cross-fold contamination
\item Manages hierarchical dependencies between preprocessing, feature 
      extraction, and model training stages
\item Validates cache hits against both parameter configurations and 
      fold assignments
\end{enumerate}

These requirements are not addressed in standard documentation, and incorrect 
implementation risks subtle bugs where models trained on Subject A are 
erroneously reused for evaluating Subject B.

This thesis provides a reference implementation demonstrating these patterns 
in a real LOSO scenario (128-fold cross-validation on Sleep-EDF), with 
explicit measurement of cache performance (hit rates, time savings) and 
validation of reproducibility. The contribution is not a new MLOps platform, 
but a validated approach for extending existing fingerprinting 
techniques to cross-validation contexts.

\chapter{Methodology}
